% chapters/ch11-state-and-time.tex
\chapter{Modelling State and Time}

\begin{learningobjectives}
In this chapter you will learn how to model concepts that evolve over time.
You will be able to distinguish entities from events, choose between storing state versus deriving it, and express time-based business constraints clearly.
You will also learn common patterns for modelling status, history, and temporal correctness.
\end{learningobjectives}

\section{Why time changes everything}
Many beginner schemas look correct until time enters the picture.
Once you need history, active versus inactive states, or “what was true at a past moment”, naive structures stop working.

Time introduces two challenges.
First, facts can change: a member changes email, a copy changes status, an appointment gets cancelled.
Second, the system must preserve history: you want to know what happened, not just what is true right now.

Relational design remains stable under time if you decide, explicitly, what should be stored as a changing attribute and what should be represented as a sequence of events.

\begin{definitionbox}
A \term{state} is a property of an entity at a point in time (e.g., copy is available).
An \term{event} is a record of something that happened at a time (e.g., a loan checkout).
\end{definitionbox}

\section{Choosing authoritative truth}
When both events and current state are stored, one must be authoritative.
If the event stream says “returned” but the state says “on loan,” the model is undefined.

In a library, the authoritative truth might be the event stream of loans.
In a clinic, the authoritative truth might be the current appointment schedule.
The choice depends on what you must audit and how often state changes.

\section{Entities versus events}
A reliable modelling decision is to treat “things” as entities and “happenings” as events.
Loans and appointments are not static links; they are happenings with time and attributes.
This is why \rel{Loan} and \rel{Appointment} are typically entities in a relational schema, even though they connect two other entities.

If you model an event as a static relationship, you often lose:
\begin{itemize}[leftmargin=*,nosep]
  \item history (you overwrite old values),
  \item constraint clarity (active state becomes unclear),
  \item and auditability (you cannot explain how you reached the current state).
\end{itemize}

\begin{examplebox}
\textbf{Bad idea: store current borrower on Copy}\par\smallskip
\[
\rel{Copy}(\underline{copy\_id}, title\_id \rightarrow \rel{Title}, current\_member\_id?)
\]
This loses history and makes “when was it borrowed?” hard to represent.
A \rel{Loan} event is a cleaner model.
\end{examplebox}

\section{Storing state versus deriving state}
A design often has two options.

Option A: store the state directly (e.g., \texttt{Copy.status = available/on\_loan/lost}).

Option B: derive the state from event history (e.g., “a copy is on loan if it has an active loan”).

Both options can be correct.
The choice depends on how the domain behaves and what needs to be enforced.

\subsection{When storing state helps}
Storing state is useful when:
\begin{itemize}[leftmargin=*,nosep]
  \item state changes are not always implied by one event stream,
  \item you need explicit human decisions (e.g., marking a copy as lost),
  \item you want fast access to current state and you can enforce consistency.
\end{itemize}

\subsection{When deriving state helps}
Deriving state is useful when:
\begin{itemize}[leftmargin=*,nosep]
  \item “current state” is fully implied by event history,
  \item you want a single source of truth (events),
  \item correctness depends on history (auditing, analytics).
\end{itemize}

In many real designs, you combine both:
events as the source of truth, with a stored current-state attribute as a cache.
If you do that, you must define consistency rules clearly, otherwise you create contradictory truths.

\begin{notebox}
If you store both events and a current-state attribute, you must decide which one is authoritative.
If both can disagree, the model is not well-defined.
\end{notebox}

\section{Status as a state machine}
Many domains have “status” attributes: appointment status, reservation status, loan status.
Status modelling is often handled poorly because it is treated as a free-form string.

At minimum, status should come from a defined set:
\[
status \in \{\text{scheduled, cancelled}\}
\]
That is a domain constraint.

For more complex domains, statuses behave like a state machine.
Some transitions are allowed; others are not.
A reservation might transition from \texttt{active} to \texttt{fulfilled} or \texttt{cancelled}, but not back to \texttt{active} without a new action.
If your book intends to teach correctness, it is useful to express this as rules.

\begin{figure}[htbp]
\centering
\includegraphics[width=0.85\textwidth]{loan-state-machine.pdf}
\caption{A status lifecycle can be modeled as explicit transitions.}
\end{figure}

\begin{examplebox}
\textbf{Status transition rules (reservation)}\\
R1: active \(\rightarrow\) fulfilled is allowed.\\
R2: active \(\rightarrow\) cancelled is allowed.\\
R3: fulfilled \(\rightarrow\) active is not allowed.
\end{examplebox}

Implementation details can wait.
The modelling point is that status is not only a value; it is behaviour.

\subsection{Diagram 1: reservation lifecycle (described)}
Imagine a simple state diagram with three nodes: \texttt{active}, \texttt{fulfilled}, \texttt{cancelled}.
Arrows go from \texttt{active} to the other two states.
There are no arrows back to \texttt{active}.

\subsection{Diagram 2: loan lifecycle (described)}
Imagine a diagram with nodes \texttt{checked\_out}, \texttt{returned}, \texttt{lost}.
Arrows go from \texttt{checked\_out} to \texttt{returned} and \texttt{lost}.
A \texttt{lost} item may later be marked \texttt{returned}, but that transition should be explicit in rules.

\section{Temporal constraints patterns}
The hardest constraints in introductory relational design are temporal.
Two common patterns appear repeatedly.

\subsection{Pattern 1: No overlap (interval conflicts)}
Clinic scheduling is a textbook example.
Each appointment defines an interval:
\[
[start\_time,\; start\_time + duration)
\]
The constraint “a doctor cannot have overlapping scheduled appointments” means:
for any doctor, no two scheduled appointment intervals may overlap.

This constraint cannot be captured by a simple uniqueness constraint when durations vary.
You can still specify it clearly as a business constraint.

\begin{examplebox}
\textbf{No-overlap constraint (conceptual)}\\
For any doctor, for any two appointments \(A\) and \(B\) with status = scheduled,\\
the intervals of \(A\) and \(B\) must not overlap.
\end{examplebox}

The schema must store enough information to evaluate overlap.
That means start time and a duration (or end time).
If you only store a start time with variable duration hidden elsewhere, you cannot validate the constraint.

\subsection{Pattern 2: One active record per entity}
In the library, the core rule is:
a copy has at most one active loan at a time.

This is a temporal constraint expressed via state.
If “active” is defined as “return\_date is missing”, then the rule becomes:
for each copy, there may be at most one loan with return\_date missing.
This is a constraint over a filtered subset of rows.

\begin{examplebox}
\textbf{One-active constraint (conceptual)}\\
For each copy, there must be at most one loan such that return\_date is missing.
\end{examplebox}

The modelling lesson is that “active” must be defined precisely.
If “active” is not defined, correctness cannot be enforced or tested.

\subsection{Pattern 3: Effective dating}
Effective dating stores validity periods explicitly:
\[
valid\_from, valid\_to
\]
The rule is that periods must not overlap for the same entity when only one value should be active.
This pattern is common for prices, memberships, and roles.

\section{History: overwrite versus append}
When values change over time, you have two general approaches.

Overwrite means you store only the current value.
This is simpler but loses history.

Append means you store changes as new rows (events or history records).
This preserves history but increases structure.

Both can be correct depending on the domain.
For teaching, it helps to show a common pattern:
a current-value table plus an optional history table.

\begin{examplebox}
\textbf{Optional history pattern}\\
Current state:
\[
\rel{Member}(\underline{member\_id}, name, email)
\]
History:
\[
\begin{aligned}
\rel{MemberEmailHistory}(&\underline{member\_id} \rightarrow \rel{Member}, \\
&\underline{changed\_at}, old\_email, new\_email)
\end{aligned}
\]
\end{examplebox}

This preserves important changes without forcing every attribute into a history structure.

\section{Common mistakes}
Time modelling mistakes are common and costly.

Mistake: Storing only current state when history is required.\\
Fix: Add event or history relations for changes that matter.

Mistake: Defining “active” informally.\\
Fix: Write a precise definition tied to attributes and time.

Mistake: Mixing end times and durations inconsistently.\\
Fix: Choose one representation and derive the other if needed.

Mistake: Allowing overlapping validity periods unintentionally.\\
Fix: Add no-overlap constraints for effective-dated records.

\section{Mini checklist}
Before you finalize time modelling, check:
\begin{itemize}[leftmargin=*,nosep]
  \item Which facts are events and which are state?
  \item Which source of truth is authoritative (events or state)?
  \item How is “active” defined for each temporal rule?
  \item Are overlap rules explicit for intervals or validity periods?
  \item Is the time representation consistent across relations?
  \item Do you need history, and if so, where is it stored?
  \item Are status transitions specified as rules?
\end{itemize}

\begin{chaptersummary}
Time forces clarity.
You must distinguish entities from events, and decide whether state is stored, derived, or both.
Status attributes should be constrained and, in richer domains, treated as state machines with explicit allowed transitions.
Temporal constraints such as “no overlap” and “one active at a time” must be defined precisely, because correctness depends on those definitions.
\end{chaptersummary}

\begin{exercisebox}
Answer in full sentences.
\begin{enumerate}[label=\arabic*.]
  \item In a clinic model, define precisely what it means for two appointments to overlap.
  \item In a library model, define precisely what “active loan” means and propose two alternative definitions.
  \item Choose one entity and decide which of its attributes should keep history. Justify why.
  \item Write three status transition rules for a concept in your chosen domain.
  \item Identify one place where storing both events and current state might help, and state which one is authoritative.
  \item Write a validity-period rule using valid\_from and valid\_to for a membership.
  \item Propose a no-overlap constraint for room bookings and define the overlap test.
  \item Give an example of a temporal rule that cannot be fully enforced and explain how to detect violations.
\end{enumerate}
\end{exercisebox}

In the next chapter we examine advanced relationship modelling.
The time and state patterns here will reappear in specializations and optionality decisions.
